% ============================================================
%  CHAPTER 18: SOUND WAVES
%  The Physics Codex: A Complete University Foundation
%  Korede Samuel Omotosho (Falcon)
%  Aurikrex Learning Systems, 2026
%
%  File: chapters/part3/ch18_sound.tex
% ============================================================

\chapter{Sound Waves}
\label{ch:sound}
\minitoc
\newpage

\chapterquote{Music is the arithmetic of sounds as
optics is the geometry of light.}{Claude Debussy}

\begin{objectivebox}
By the end of this chapter, you should be able to:
\begin{enumerate}[noitemsep]
  \item Describe the nature of sound as a longitudinal
  mechanical wave
  \item Explain how sound is produced and detected
  \item Calculate the speed of sound in different
  media and apply it to problems
  \item Define and calculate frequency, pitch, loudness
  and quality of sound
  \item Explain and calculate the Doppler effect
  \item Describe echoes and use them to calculate
  distances
  \item Explain resonance in sound and give practical
  examples
  \item Define the decibel scale and perform basic
  calculations
  \item Distinguish between ultrasound and infrasound
  and describe their applications
\end{enumerate}
\end{objectivebox}

% ============================================================
\section{Intuition: The Nature of Sound}
% ============================================================

Strike a tuning fork and hold it near your ear. The
sound you hear is produced by the vibrating prongs of
the fork. Each prong pushes the air molecules in front
of it, creating a region of higher pressure (a
\textbf{compression}). As the prong swings back, it
leaves a region of lower pressure (a
\textbf{rarefaction}). These pressure oscillations
travel outward through the air at approximately
$340\ \text{m\,s}^{-1}$, reach your eardrum, cause it
to vibrate at the same frequency, and your brain
interprets this as sound.

Sound is a \textbf{longitudinal mechanical wave} of
pressure variations travelling through a medium. Three
things are absolutely essential:
\begin{enumerate}[noitemsep]
  \item A vibrating source
  \item A material medium to carry the wave
  \item A detector (ear, microphone, instrument)
\end{enumerate}

Without a medium, there is no sound. The famous
tagline from the film \textit{Alien} --- ``In space,
no one can hear you scream'' --- is physically accurate.
Sound cannot travel through a vacuum.

% ============================================================
\section{Production and Propagation of Sound}
% ============================================================

\begin{diagbox}[Tuning Fork Producing Sound Waves]
\begin{center}
\begin{tikzpicture}[scale=1.0, font=\small]

  % Tuning fork
  \fill[gray!60] (0,3.5) rectangle (0.4,6.5);
  \fill[gray!60] (0,5.5) .. controls (-0.5,5.5) and
    (-0.5,7.5) .. (0,7.5) -- (0.3,7.5) ..
    controls (-0.2,7.5) and (-0.2,5.5) .. (0.3,5.5);
  \fill[gray!60] (0.4,5.5) .. controls (0.9,5.5)
    and (0.9,7.5) .. (0.4,7.5) -- (0.7,7.5) ..
    controls (1.2,7.5) and (1.2,5.5) .. (0.7,5.5);
  \fill[gray!60] (0.3,3.5) rectangle (0.7,5.5);

  % Prong vibration arrow
  \draw[->, thick, codexred] (-0.3,7.0) -- (-0.7,7.0);
  \draw[->, thick, codexblue] (1.0,7.0) -- (1.4,7.0);
  \node[codexred, font=\tiny] at (-0.5,7.3)
    {compress};
  \node[codexblue, font=\tiny] at (1.2,7.3)
    {compress};

  % Compressions (C) and rarefactions (R) spreading right
  % Compression 1
  \foreach \x in {2.0, 2.1, 2.2} {
    \fill[codexblue!60] (\x, 4.5) circle (3pt);
    \fill[codexblue!60] (\x, 5.0) circle (3pt);
    \fill[codexblue!60] (\x, 5.5) circle (3pt);
    \fill[codexblue!60] (\x, 6.0) circle (3pt);
    \fill[codexblue!60] (\x, 6.5) circle (3pt);
  }
  \node[codexblue, font=\tiny] at (2.1,4.1) {C};

  % Rarefaction 1 (sparse)
  \foreach \x in {3.0, 3.5} {
    \fill[codexblue!30] (\x, 4.5) circle (2pt);
    \fill[codexblue!30] (\x, 5.5) circle (2pt);
    \fill[codexblue!30] (\x, 6.5) circle (2pt);
  }
  \node[gray, font=\tiny] at (3.25,4.1) {R};

  % Compression 2
  \foreach \x in {4.2, 4.3, 4.4} {
    \fill[codexblue!60] (\x, 4.5) circle (3pt);
    \fill[codexblue!60] (\x, 5.0) circle (3pt);
    \fill[codexblue!60] (\x, 5.5) circle (3pt);
    \fill[codexblue!60] (\x, 6.0) circle (3pt);
    \fill[codexblue!60] (\x, 6.5) circle (3pt);
  }
  \node[codexblue, font=\tiny] at (4.3,4.1) {C};

  % Rarefaction 2
  \foreach \x in {5.2, 5.7} {
    \fill[codexblue!30] (\x, 4.5) circle (2pt);
    \fill[codexblue!30] (\x, 5.5) circle (2pt);
    \fill[codexblue!30] (\x, 6.5) circle (2pt);
  }
  \node[gray, font=\tiny] at (5.45,4.1) {R};

  % Compression 3
  \foreach \x in {6.4, 6.5, 6.6} {
    \fill[codexblue!60] (\x, 4.5) circle (3pt);
    \fill[codexblue!60] (\x, 5.0) circle (3pt);
    \fill[codexblue!60] (\x, 5.5) circle (3pt);
    \fill[codexblue!60] (\x, 6.0) circle (3pt);
    \fill[codexblue!60] (\x, 6.5) circle (3pt);
  }
  \node[codexblue, font=\tiny] at (6.5,4.1) {C};

  % Wave direction arrow
  \draw[->, very thick, codexgold]
    (2.0,3.7) -- (7.5,3.7)
    node[right]{direction of travel};

  % Pressure graph below
  \begin{scope}[shift={(1.5,0)}]
  \begin{axis}[
    at={(0pt,-30pt)},
    width=9cm, height=3.5cm,
    xlabel={Distance from source},
    ylabel={Pressure},
    xmin=0, xmax=6,
    ymin=-1.3, ymax=1.3,
    grid=major, grid style={gray!15},
    xtick=\empty,
    ytick={-1,0,1},
    yticklabels={low $P$,$P_0$,high $P$},
  ]
    \addplot[codexblue, very thick, domain=0:6,
      samples=100]
      {sin(deg(x*3.14159))};
    \node[font=\tiny, codexblue] at (axis cs:0.5,0.8)
      {C};
    \node[font=\tiny, gray] at (axis cs:1.5,-0.8)
      {R};
    \node[font=\tiny, codexblue] at (axis cs:2.5,0.8)
      {C};
    \node[font=\tiny, gray] at (axis cs:3.5,-0.8)
      {R};
    \node[font=\tiny, codexblue] at (axis cs:4.5,0.8)
      {C};
    \node[font=\tiny, gray] at (axis cs:5.5,-0.8)
      {R};
  \end{axis}
  \end{scope}

\end{tikzpicture}
\end{center}
\small Sound is a longitudinal wave: the air molecules
oscillate back and forth (parallel to the direction
of travel), creating alternating compressions (high
pressure) and rarefactions (low pressure) that
propagate outward.
\end{diagbox}

% ============================================================
\section{Speed of Sound}
% ============================================================

Sound travels at different speeds through different
materials. In general, sound travels fastest through
solids (tightly packed molecules transmit vibrations
quickly), slower through liquids, and slowest through
gases (molecules are far apart and collide less
frequently).

\begin{table}[H]
\centering
\caption{Speed of Sound in Various Media}
\label{tab:sound_speed}
\begin{tabular}{ll}
\toprule
\textbf{Medium} & \textbf{Speed of Sound (m\,s$^{-1}$)}\\
\midrule
Air at 0°C & 331\\
Air at 20°C & 343\\
Air at 30°C & 349\\
Water (fresh, 20°C) & 1480\\
Seawater & 1520\\
Wood (oak) & 3850\\
Concrete & 3400\\
Steel & 5100\\
Glass & 5000\\
Aluminium & 6320\\
\bottomrule
\end{tabular}
\end{table}

\begin{keybox}[Speed of Sound in Air and Temperature]
The speed of sound in air increases with temperature:
\[
  v \approx 331 + 0.6\theta\ \text{m\,s}^{-1}
\]
where $\theta$ is the temperature in °C. At $20°\text{C}$:
$v \approx 331 + 0.6(20) = 343\ \text{m\,s}^{-1}$

The speed of sound is \textbf{independent of frequency}
and \textbf{independent of pressure} (for ideal gases).
It depends only on the properties of the medium
(temperature, density, and elasticity).
\end{keybox}

\begin{examplebox}[18.1]
\stepcounter{workedexample}
During a thunderstorm, a student sees a lightning
flash and hears the thunder $4.5\ \text{s}$ later.
The temperature is $30°\text{C}$.\\
(a) Calculate the speed of sound.\\
(b) Calculate the distance to the lightning strike.\\
(c) Why does the student see the lightning before
hearing the thunder?
\end{examplebox}

\begin{solutionbox}
\textbf{(a)} Speed of sound at $30°\text{C}$:
\[
  v = 331 + 0.6(30) = 331 + 18 = 349\ \text{m\,s}^{-1}
\]

\textbf{(b)} Distance:
\[
  d = v \times t = 349 \times 4.5
  = 1570.5\ \text{m} \approx 1.57\ \text{km}
\]

\textbf{(c)} Light travels at $3\times10^8\ \text{m\,s}^{-1}$,
approximately $10^6$ times faster than sound. The light from
the lightning reaches the student almost instantaneously
($\sim 5\times10^{-6}\ \text{s}$), while the sound
takes $4.5\ \text{s}$. This is why you always see
lightning before hearing thunder.

\begin{realworldbox}[Counting Seconds for Safety]
The $3$-second rule: count the seconds between seeing
lightning and hearing thunder, then divide by $3$ to
get the approximate distance in kilometres. At
$3\ \text{seconds}$, the storm is about $1\ \text{km}$
away. At $3\ \text{seconds}$ or less, lightning is close
enough to be dangerous --- shelter immediately. This
simple rule uses $v_{sound} \approx 333\ \text{m\,s}^{-1}
\approx 1\ \text{km per 3 seconds}$. It works well in
tropical temperatures like those across Nigeria.
\end{realworldbox}
\end{solutionbox}

% ============================================================
\section{Characteristics of Sound}
% ============================================================

\begin{defbox}[Pitch, Loudness and Quality]
\textbf{Pitch} is the perceived highness or lowness
of a sound. It is determined primarily by the
\textbf{frequency} of the sound wave.
Higher frequency $\implies$ higher pitch.
The human ear can detect frequencies from about
$20\ \text{Hz}$ to $20\,000\ \text{Hz}$.

\textbf{Loudness} (volume) is the perceived intensity
of a sound. It is determined by the \textbf{amplitude}
of the sound wave. Greater amplitude $\implies$ louder
sound. Loudness is measured in decibels (dB).

\textbf{Quality} (timbre) is what makes a trumpet
sound different from a violin even when both play the
same note at the same loudness. It is determined by
the \textbf{waveform} of the sound --- specifically,
the mixture of harmonics (overtones) present in the
sound. A pure tone (single frequency) has a sinusoidal
waveform. A rich musical tone contains many harmonics.
\end{defbox}

\begin{diagbox}[Pitch, Loudness and Quality as Waveforms]
\begin{center}
\begin{tikzpicture}
  % High pitch vs low pitch
  \begin{axis}[
    name=p1,
    width=11cm, height=3cm,
    title={Pitch: High frequency = high pitch},
    xlabel={}, ylabel={},
    xmin=0, xmax=4*pi,
    ymin=-1.5, ymax=1.5,
    xtick=\empty, ytick=\empty,
    grid=major, grid style={gray!10},
    legend style={at={(0.5,-0.22)}, anchor=north, font=\tiny, /tikz/column 2/.style={column sep=3pt}},
  ]
    \addplot[codexblue, very thick, domain=0:4*pi,
      samples=200]{sin(deg(4*x))};
    \addlegendentry{High frequency (high pitch)};
    \addplot[codexred, thick, dashed, domain=0:4*pi,
      samples=100]{sin(deg(1.5*x))};
    \addlegendentry{Low frequency (low pitch)};
  \end{axis}

  % Loud vs quiet
  \begin{axis}[
    name=p2,
    at={(p1.below south west)},
    anchor=above north west,
    yshift=-0.1cm,
    width=11cm, height=3cm,
    title={Loudness: Large amplitude = loud sound},
    xlabel={}, ylabel={},
    xmin=0, xmax=4*pi,
    ymin=-1.5, ymax=1.5,
    xtick=\empty, ytick=\empty,
    grid=major, grid style={gray!10},
    legend style={at={(0.5,-0.22)}, anchor=north, font=\tiny, /tikz/column 2/.style={column sep=3pt}},
  ]
    \addplot[codexblue, very thick, domain=0:4*pi,
      samples=200]{1.2*sin(deg(2*x))};
    \addlegendentry{Large amplitude (loud)};
    \addplot[codexred, thick, dashed, domain=0:4*pi,
      samples=200]{0.4*sin(deg(2*x))};
    \addlegendentry{Small amplitude (quiet)};
  \end{axis}

  % Quality (timbre) - pure vs complex
  \begin{axis}[
    name=p3,
    at={(p2.below south west)},
    anchor=above north west,
    yshift=-0.1cm,
    width=11cm, height=3cm,
    title={Quality: Different waveforms, same frequency},
    xlabel={}, ylabel={},
    xmin=0, xmax=4*pi,
    ymin=-1.5, ymax=1.5,
    xtick=\empty, ytick=\empty,
    grid=major, grid style={gray!10},
    legend style={at={(0.5,-0.22)}, anchor=north, font=\tiny, /tikz/column 2/.style={column sep=3pt}},
  ]
    \addplot[codexblue, very thick, domain=0:4*pi,
      samples=200]{sin(deg(2*x))};
    \addlegendentry{Pure tone (sine wave)};
    \addplot[codexred, thick, dashed, domain=0:4*pi,
      samples=200]
      {0.7*sin(deg(2*x)) + 0.3*sin(deg(4*x))
      + 0.15*sin(deg(6*x))};
    \addlegendentry{Rich tone (harmonics added)};
    \node[font=\tiny, anchor=south east, fill=white, inner sep=1pt]
      at (axis cs:12,-1.35) {Time};
  \end{axis}
\end{tikzpicture}
\end{center}
\end{diagbox}

% ============================================================
\section{The Decibel Scale}
% ============================================================

\begin{definitionbox}[The Decibel Scale]
The human ear responds to sound intensity over an
enormous range --- from the threshold of hearing
($I_0 = 10^{-12}\ \text{W\,m}^{-2}$) to the threshold
of pain ($\approx 1\ \text{W\,m}^{-2}$), a ratio of
$10^{12}$. To handle this, sound level $L$ is measured
in \textbf{decibels (dB)}:
\[
  L = 10\log_{10}\left(\frac{I}{I_0}\right)\ \text{dB}
\]
where $I$ is the intensity and $I_0 = 10^{-12}\ \text{W\,m}^{-2}$
is the reference (threshold of hearing).

Key values:
\begin{itemize}[noitemsep]
  \item Every \textbf{10 dB increase} corresponds to
  a \textbf{10-fold increase in intensity}
  \item Every \textbf{3 dB increase} corresponds to
  approximately a \textbf{doubling of intensity}
\end{itemize}
\end{definitionbox}

\begin{table}[H]
\centering
\caption{Typical Sound Levels}
\label{tab:sound_levels}
\begin{tabular}{lll}
\toprule
\textbf{Sound} & \textbf{Level (dB)} &
\textbf{Intensity (W\,m$^{-2}$)}\\
\midrule
Threshold of hearing & 0 & $10^{-12}$\\
Rustling leaves & 10 & $10^{-11}$\\
Whisper & 30 & $10^{-9}$\\
Quiet conversation & 50 & $10^{-7}$\\
Normal conversation & 60 & $10^{-6}$\\
Lagos traffic & 80 & $10^{-4}$\\
Generator (nearby) & 90 & $10^{-3}$\\
Live concert & 110 & $10^{-1}$\\
Threshold of pain & 120--130 & $1$--$10$\\
Jet engine (nearby) & 140 & $100$\\
\bottomrule
\end{tabular}
\end{table}

\begin{examplebox}[18.2]
\stepcounter{workedexample}
A generator produces a sound level of $90\ \text{dB}$.
A second identical generator is switched on next to it.\\
(a) What is the new sound level?\\
(b) A third identical generator is added. What is
the level now?\\
(c) How many generators would be needed to raise the
level to $100\ \text{dB}$?
\end{examplebox}

\begin{solutionbox}
\textbf{(a)} Each generator has intensity
$I = 10^{-3}\ \text{W\,m}^{-2}$ (from Table).

Two generators: $I_{total} = 2 \times 10^{-3}\ \text{W\,m}^{-2}$
\[
  L = 10\log_{10}\left(\frac{2\times10^{-3}}{10^{-12}}\right)
  = 10\log_{10}(2\times10^9)
  = 10(9 + \log_{10}2)
  = 10(9 + 0.301) = 93.0\ \text{dB}
\]

\textbf{(b)} Three generators: $I = 3\times10^{-3}$
\[
  L = 10\log_{10}(3\times10^9) = 10(9+0.477)
  = 94.8\ \text{dB}
\]

\textbf{(c)} For $100\ \text{dB}$, need intensity
$I = 10^{-2}\ \text{W\,m}^{-2}$:
\[
  n = \frac{10^{-2}}{10^{-3}} = 10\ \text{generators}
\]
You need \textit{ten} generators to raise the level
by just $10\ \text{dB}$. This illustrates why simply
adding more generators is an inefficient way to
increase loudness.
\end{solutionbox}

% ============================================================
\section{Echoes and Reverberation}
% ============================================================

\begin{definitionbox}[Echo and Reverberation]
An \textbf{echo} is a reflected sound that is heard
distinctly as a separate sound from the original.
For an echo to be perceived as separate, the reflected
sound must arrive at least $\frac{1}{10}\ \text{s}$
after the direct sound (the ear's persistence of
sound is about $0.1\ \text{s}$).

Minimum distance for an echo:
\[
  d_{min} = \frac{v \times t_{min}}{2}
  = \frac{340 \times 0.1}{2} = 17\ \text{m}
\]

\textbf{Reverberation} is the persistence of sound
in an enclosed space due to multiple reflections.
Unlike an echo (which has a clear time separation),
reverberation produces a gradual decay of sound.
Desirable in concert halls; undesirable in classrooms
(where it reduces speech clarity).
\end{definitionbox}

\begin{examplebox}[18.3]
\stepcounter{workedexample}
A ship uses an echo sounder (SONAR) to measure
the depth of the ocean. An ultrasound pulse is
transmitted downward and the echo is received
$0.24\ \text{s}$ later. The speed of sound in
seawater is $1520\ \text{m\,s}^{-1}$.\\
(a) Calculate the depth of the ocean at this point.\\
(b) The ship moves and the echo now takes $0.48\ \text{s}$.
Has the ocean floor sloped up or down? By how much?
\end{examplebox}

\begin{solutionbox}
\textbf{(a)} The sound travels down to the ocean
floor and back (total distance $= 2d$):
\[
  d = \frac{v \times t}{2}
  = \frac{1520 \times 0.24}{2}
  = \frac{364.8}{2} = 182.4\ \text{m}
\]

\textbf{(b)} New depth:
\[
  d' = \frac{1520 \times 0.48}{2}
  = \frac{729.6}{2} = 364.8\ \text{m}
\]

The ocean floor has sloped \textbf{down} by:
$364.8 - 182.4 = 182.4\ \text{m}$
\end{solutionbox}

% ============================================================
\section{The Doppler Effect}
% ============================================================

Stand near a road and listen as an ambulance approaches
with its siren on. The pitch of the siren is noticeably
higher as the ambulance approaches than it is after
it has passed. The frequency of the sound \textit{heard}
by a listener is different from the frequency emitted
by the source when there is relative motion between
the source and the listener. This is the
\textbf{Doppler effect}.

\begin{keybox}[Intuition: Why the Doppler Effect Occurs]
Imagine a source emitting $f$ waves per second. As it
moves toward you, it is chasing the wavefronts it has
already emitted --- they bunch up in front of it,
reducing the wavelength you receive. Shorter wavelength
at the same wave speed means higher frequency --- higher
pitch.

As it moves away, the source stretches the wavefronts
behind it, increasing the wavelength you receive.
Longer wavelength means lower frequency --- lower pitch.
\end{keybox}

\begin{diagbox}[Doppler Effect: Moving Source]
\begin{center}
\begin{tikzpicture}[scale=0.8, font=\small]
  % One-dimensional cross-section of wavefronts at a common instant.
  % The source moves right, so successive fronts are closer ahead and
  % farther apart behind it.
  \fill[codexgreen] (4.6,0) circle (5pt);
  \node[above, codexgreen, align=center] at (4.6,0.2)
    {Observer A\\ahead; higher pitch};

  \fill[codexred] (-5.4,0) circle (5pt);
  \node[below, codexred, align=center] at (-5.4,-0.2)
    {Observer B\\behind; lower pitch};

  % Compressed wavefront spacing in front of the source.
  \foreach \x in {0.7,1.4,2.1} {
    \draw[codexblue, thick] (\x,-0.65) -- (\x,0.65);
  }
  % Expanded wavefront spacing behind the source.
  \foreach \x in {-1.4,-2.8,-4.2} {
    \draw[codexblue, thick] (\x,-0.65) -- (\x,0.65);
  }

  % Current source position and direction of travel.
  \fill[codexblue] (0,0) circle (6pt);
  \node[above, codexblue] at (0,0.75) {Source};
  \draw[->, very thick, codexblue] (0,1.35) -- (1.0,1.35)
    node[right, font=\small]{$v_s$};

  % Wavelengths are the separations between adjacent fronts.
  \draw[<->, codexgreen, thick] (0.7,-1.0) -- (1.4,-1.0);
  \node[below, codexgreen, font=\tiny] at (1.05,-1.0)
    {$\lambda_{\mathrm{front}}$ (short)};
  \draw[<->, codexred, thick] (-2.8,-1.0) -- (-1.4,-1.0);
  \node[below, codexred, font=\tiny] at (-2.1,-1.0)
    {$\lambda_{\mathrm{rear}}$ (long)};

  \node[font=\tiny, gray, align=center] at (0,-2.0)
    {Schematic wavefront spacing at one instant};
\end{tikzpicture}
\end{center}
\end{diagbox}

\begin{lawbox}[Doppler Effect Formula]
For a source of frequency $f_0$ moving at speed $v_s$
and an observer moving at speed $v_o$, with the speed
of sound $v$:

\[
  f_{observed} = f_0 \times
  \frac{v \pm v_o}{v \mp v_s}
\]

\textbf{Sign convention:}
\begin{itemize}[noitemsep]
  \item Use $+v_o$ (numerator) when observer moves
  \textit{toward} source
  \item Use $-v_o$ (numerator) when observer moves
  \textit{away from} source
  \item Use $-v_s$ (denominator) when source moves
  \textit{toward} observer
  \item Use $+v_s$ (denominator) when source moves
  \textit{away from} observer
\end{itemize}

For a \textbf{stationary observer} and \textbf{moving
source}:
\[
  f_{obs} = f_0 \times \frac{v}{v \mp v_s}
\]
(use $-$ when source approaches, $+$ when it recedes)
\end{lawbox}

\begin{examplebox}[18.4]
\stepcounter{workedexample}
An ambulance travelling at $30\ \text{m\,s}^{-1}$
sounds its siren at $800\ \text{Hz}$. The speed of
sound is $340\ \text{m\,s}^{-1}$.\\
(a) What frequency does a stationary pedestrian
hear as the ambulance approaches?\\
(b) What frequency does the pedestrian hear as the
ambulance moves away?\\
(c) What is the change in frequency as the
ambulance passes?
\end{examplebox}

\begin{solutionbox}
$f_0 = 800\ \text{Hz}$, $v_s = 30\ \text{m\,s}^{-1}$,
$v = 340\ \text{m\,s}^{-1}$, $v_o = 0$

\textbf{(a)} Ambulance approaching
(use $-v_s$ in denominator):
\[
  f_{approach} = 800 \times \frac{340}{340 - 30}
  = 800 \times \frac{340}{310}
  = 800 \times 1.0968
  = 877.4\ \text{Hz} \approx 877\ \text{Hz}
\]

\textbf{(b)} Ambulance receding
(use $+v_s$ in denominator):
\[
  f_{recede} = 800 \times \frac{340}{340 + 30}
  = 800 \times \frac{340}{370}
  = 800 \times 0.9189
  = 735.1\ \text{Hz} \approx 735\ \text{Hz}
\]

\textbf{(c)} Change in frequency:
\[
  \Delta f = 877 - 735 = 142\ \text{Hz}
\]

The pedestrian hears a dramatic drop of $142\ \text{Hz}$
as the ambulance passes --- from a noticeably high pitch
to a noticeably low pitch. This is the classic sound
every Nigerian (and everyone else) recognises when
an emergency vehicle passes at speed.
\end{solutionbox}

\begin{realworldbox}[Doppler in Medicine and Astronomy]
The Doppler effect is used in \textbf{Doppler
ultrasound} to measure blood flow in medical
diagnostics. Ultrasound pulses reflect off moving
red blood cells; the frequency shift reveals the
speed and direction of blood flow. Cardiologists use
this to detect blocked arteries and abnormal heart
valve function --- entirely non-invasive, no surgery
required.

In astronomy, the \textbf{redshift} of light from
distant galaxies (a Doppler shift toward lower
frequencies as galaxies recede) was the evidence
that first convinced scientists the universe is
expanding. Edwin Hubble measured this in 1929. Every
galaxy is moving away from us, and the more distant
it is, the faster it recedes. The Doppler effect
connects a tuning fork on a road to the expansion
of the universe.
\end{realworldbox}

% ============================================================
\section{Resonance in Sound}
% ============================================================

When a sound wave of the right frequency is applied to
an air column or string, the system vibrates with very
large amplitude --- it \textbf{resonates}. This occurs
at the natural frequencies (harmonics) of the system.

\begin{examplebox}[18.5]
\stepcounter{workedexample}
A resonance tube (closed at one end) is found to
resonate at $17\ \text{cm}$ and $51\ \text{cm}$
when a tuning fork of frequency $500\ \text{Hz}$
is held over it. The speed of sound is
$340\ \text{m\,s}^{-1}$.\\
(a) Explain why resonance occurs at these two lengths.\\
(b) Calculate the wavelength from the resonance
lengths.\\
(c) Calculate the speed of sound and compare with
the given value.\\
(d) Calculate the end correction.
\end{examplebox}

\begin{solutionbox}
\textbf{(a)} A closed pipe has a node at the closed
end and an antinode at the open end. The first
resonance occurs at $L_1 = \lambda/4$ and the
second at $L_2 = 3\lambda/4$ (1st and 3rd harmonics
of the closed pipe).

\textbf{(b)} From $L_2 - L_1 = \lambda/2$:
\[
  \lambda = 2(L_2 - L_1) = 2(51 - 17)
  = 2 \times 34 = 68\ \text{cm} = 0.68\ \text{m}
\]

\textbf{(c)} Speed of sound:
\[
  v = f\lambda = 500 \times 0.68
  = 340\ \text{m\,s}^{-1}
\]
Exact agreement. The calculation is self-consistent.

\textbf{(d)} For a closed pipe with end correction $e$:
\[
  L_1 + e = \frac{\lambda}{4}
  \implies e = \frac{\lambda}{4} - L_1
  = \frac{0.68}{4} - 0.17 = 0.17 - 0.17 = 0\ \text{m}
\]

In this idealised problem, the end correction is zero.
In real experiments, the antinode is slightly beyond
the open end, giving $e \approx 0.3r$ where $r$ is
the tube radius.
\end{solutionbox}

% ============================================================
\section{Ultrasound and Infrasound}
% ============================================================

\begin{defbox}[Ultrasound and Infrasound]
\textbf{Audible sound:} $20\ \text{Hz}$
to $20\,000\ \text{Hz}$ (20 kHz)

\textbf{Infrasound:} frequency $< 20\ \text{Hz}$.
Produced by earthquakes, volcanoes, large animals
(elephants, whales). Travels very long distances
with little attenuation. Used by elephants to
communicate over many kilometres.

\textbf{Ultrasound:} frequency $> 20\,000\ \text{Hz}$.
Produced by bats, dolphins, dogs.
Applications:
\begin{itemize}[noitemsep]
  \item Medical imaging (foetal scans, organ imaging)
  \item SONAR (depth measurement, submarine detection)
  \item Industrial flaw detection (cracks in metals)
  \item Cleaning delicate instruments
  \item Measuring speed of blood flow (Doppler
  ultrasound)
\end{itemize}
\textbf{Advantage of ultrasound for imaging:} short
wavelength means high resolution (features smaller
than the wavelength cannot be resolved).
\end{defbox}

\begin{diagbox}[Medical Ultrasound Imaging: Principle]
\begin{center}
\begin{tikzpicture}[scale=0.95, font=\small]

  % Transducer
  \fill[gray!60] (0,3.5) rectangle (1.5,4.5);
  \draw[thick] (0,3.5) rectangle (1.5,4.5);
  \node[above, font=\tiny] at (0.75,4.5)
    {Transducer};

  % Tissue layers
  \fill[codexblue!10] (1.5,0) rectangle (8,5);
  \draw[thick] (1.5,0) -- (1.5,5);

  % Boundary 1 (e.g. skin-muscle interface)
  \fill[gray!30] (3.0,0) rectangle (3.3,5);
  \node[font=\tiny, rotate=90] at (3.15,2.5) {skin/muscle};

  % Boundary 2 (e.g. organ surface)
  \fill[codexgold!30] (5.5,0) rectangle (5.8,5);
  \node[font=\tiny, rotate=90] at (5.65,2.5) {organ wall};

  % Ultrasound pulse
  \draw[->, very thick, codexblue]
    (1.5,4.0) -- (3.0,4.0)
    node[above, midway, font=\tiny, fill=white, inner sep=1pt]{pulse};

  % Reflection 1
  \draw[->, thick, codexred, dashed]
    (3.0,3.8) -- (1.6,3.8)
    node[below, midway, font=\tiny, fill=white, inner sep=1pt]{echo 1};

  % Transmitted pulse continues
  \draw[->, thick, codexblue]
    (3.3,4.0) -- (5.5,4.0);

  % Reflection 2
  \draw[->, thick, codexred, dashed]
    (5.5,3.2) -- (1.6,3.2)
    node[below, midway, font=\tiny, fill=white, inner sep=1pt]{echo 2};

  % Time of flight labels
  \draw[<->, gray] (1.6,1.5) -- (3.0,1.5);
  \node[below, gray, font=\tiny] at (2.3,1.4) {$t_1$};
  \draw[<->, gray] (1.6,0.8) -- (5.5,0.8);
  \node[below, gray, font=\tiny] at (3.55,0.7) {$t_2$};

  % Depth formula
  \node[font=\small, align=center, fill=white, inner sep=2pt] at (7.0,1.8)
    {$d_1 = \frac{v \times t_1}{2}$\\[4pt]
    $d_2 = \frac{v \times t_2}{2}$};

  % A-scan display: brief echo pulses at increasing return times.
  \fill[black] (9.0,0) rectangle (11.0,5);
  \node[white, font=\tiny, align=center] at (10,4.72)
    {A-scan display};
  \node[white, font=\tiny, align=center] at (10,4.34)
    {amplitude vs\\time};
  \draw[->, white!80] (9.28,0.55) -- (9.28,3.65);
  \node[white, font=\tiny, rotate=90] at (9.13,2.15) {amplitude};
  \draw[->, white!80] (9.28,0.55) -- (10.82,0.55);
  \node[white, font=\tiny, anchor=west] at (10.20,0.32) {time};
  \draw[white!40, thin] (9.30,1.05) -- (10.82,1.05);
  \draw[codexgreen, very thick]
    (9.30,1.05) -- (9.62,1.05) -- (9.70,2.75) -- (9.78,1.05)
    -- (10.27,1.05) -- (10.35,2.20) -- (10.43,1.05) -- (10.82,1.05);
  \node[codexgreen, font=\tiny, anchor=south] at (9.70,2.82) {echo 1};
  \node[codexgreen, font=\tiny, anchor=south] at (10.35,2.27) {echo 2};

\end{tikzpicture}
\end{center}
\small The transducer emits a pulse and then listens
for echoes. The time delay of each echo reveals the
depth of the reflecting boundary: $d = vt/2$. Building
up many such measurements creates a cross-sectional
image of internal organs.
\end{diagbox}

% ============================================================
\section{Chapter Summary}
% ============================================================

\begin{summarybox}
\textbf{Sound:} longitudinal mechanical wave (pressure
variations); cannot travel in a vacuum.

\textbf{Speed:} $v_{air} \approx 331 + 0.6\theta$
m\,s$^{-1}$; faster in denser/stiffer media.
Fastest in solids, slowest in gases.

\textbf{Pitch:} determined by frequency.
\textbf{Loudness:} determined by amplitude.
\textbf{Quality (timbre):} determined by harmonic
content (waveform).

\textbf{Decibel:} $L = 10\log_{10}(I/I_0)$ dB;
$I_0 = 10^{-12}\ \text{W\,m}^{-2}$.
$+10\ \text{dB} \equiv \times 10$ intensity.

\textbf{Echo:} reflected sound; $d = vt/2$.
Distinct echo requires $t > 0.1\ \text{s}$,
$d > 17\ \text{m}$.

\textbf{Doppler effect:}
$f_{obs} = f_0 \frac{v \pm v_o}{v \mp v_s}$

Source approaching: $f_{obs} > f_0$.
Source receding: $f_{obs} < f_0$.

\textbf{Resonance:} maximum amplitude at natural
frequency.

\textbf{Ultrasound} ($>20\ \text{kHz}$): medical
imaging, SONAR, flaw detection.
Short wavelength $\implies$ high resolution.

\textbf{Infrasound} ($<20\ \text{Hz}$): earthquakes,
large animals; travels very long distances.
\end{summarybox}

% ============================================================
\newpage
\begin{exercisebox}[18]

\subsection*{Section A: Multiple Choice}

\textbf{A1.} Sound travels fastest through:

\mcoption{A}{Air at 0°C}
\mcoption{B}{Water}
\mcoption{C}{Steel}
\mcoption{D}{A vacuum}
\ansref{Ch18-A1}

\vspace{0.4cm}

\textbf{A2.} The pitch of a musical note is
determined primarily by its:

\mcoption{A}{Amplitude}
\mcoption{B}{Speed}
\mcoption{C}{Frequency}
\mcoption{D}{Wavelength in the instrument}
\ansref{Ch18-A2}

\vspace{0.4cm}

\textbf{A3.} A sound of intensity $10^{-8}\ \text{W\,m}^{-2}$
has a sound level of
($I_0 = 10^{-12}\ \text{W\,m}^{-2}$):

\mcoption{A}{$4\ \text{dB}$}
\mcoption{B}{$8\ \text{dB}$}
\mcoption{C}{$40\ \text{dB}$}
\mcoption{D}{$80\ \text{dB}$}
\ansref{Ch18-A3}

\vspace{0.4cm}

\textbf{A4.} A fire engine siren emits $600\ \text{Hz}$.
As it approaches at $40\ \text{m\,s}^{-1}$, a
stationary observer hears a frequency of
($v_{sound} = 340\ \text{m\,s}^{-1}$):

\mcoption{A}{$526\ \text{Hz}$}
\mcoption{B}{$600\ \text{Hz}$}
\mcoption{C}{$680\ \text{Hz}$}
\mcoption{D}{$740\ \text{Hz}$}
\ansref{Ch18-A4}

\vspace{0.4cm}

\textbf{A5.} An echo is heard $0.6\ \text{s}$ after
a clap. The reflecting wall is
($v = 340\ \text{m\,s}^{-1}$):

\mcoption{A}{$34\ \text{m}$ away}
\mcoption{B}{$51\ \text{m}$ away}
\mcoption{C}{$102\ \text{m}$ away}
\mcoption{D}{$204\ \text{m}$ away}
\ansref{Ch18-A5}

\vspace{0.4cm}

\textbf{A6.} Which of the following correctly
describes ultrasound?

\mcoption{A}{Sound above the speed of sound}
\mcoption{B}{Sound with frequency below 20 Hz}
\mcoption{C}{Sound with frequency above 20 kHz}
\mcoption{D}{Very loud sound above 120 dB}
\ansref{Ch18-A6}

\vspace{0.4cm}

\textbf{A7.} In the Doppler effect, if a source
of sound is moving away from a stationary observer:

\mcoption{A}{The observed frequency is higher than
the emitted frequency}
\mcoption{B}{The observed frequency equals the
emitted frequency}
\mcoption{C}{The observed frequency is lower than
the emitted frequency}
\mcoption{D}{The speed of sound increases}
\ansref{Ch18-A7}

\vspace{0.4cm}

\textbf{A8.} Sound cannot travel through a vacuum
because:

\mcoption{A}{There is no gravity in a vacuum}
\mcoption{B}{There are no molecules to vibrate
and transmit the pressure wave}
\mcoption{C}{The temperature is too low}
\mcoption{D}{Electromagnetic waves cancel sound
waves in a vacuum}
\ansref{Ch18-A8}

\vspace{0.4cm}

\textbf{A9.} Ten identical sound sources together
produce $70\ \text{dB}$. One source alone produces:

\mcoption{A}{$7\ \text{dB}$}
\mcoption{B}{$60\ \text{dB}$}
\mcoption{C}{$63\ \text{dB}$}
\mcoption{D}{$700\ \text{dB}$}
\ansref{Ch18-A9}

\vspace{0.4cm}

\textbf{A10.} In a resonance tube experiment (closed
pipe), the first and second resonance positions are
at $L_1$ and $L_2$. The wavelength of the sound is:

\mcoption{A}{$L_1$}
\mcoption{B}{$2L_1$}
\mcoption{C}{$L_2 - L_1$}
\mcoption{D}{$2(L_2 - L_1)$}
\ansref{Ch18-A10}

\vspace{0.6cm}

\subsection*{Section B: Theory and Calculations}

\textbf{B1.} A student stands $85\ \text{m}$ from
a large wall and claps her hands. She then claps
rhythmically so that each clap coincides with the
echo of the previous clap. The temperature is
$25°\text{C}$.\\[4pt]
(a) Calculate the speed of sound at $25°\text{C}$.\\
(b) Calculate the time between successive claps.\\
(c) Calculate the number of claps per second
(the frequency of clapping).\\
(d) She moves $10\ \text{m}$ closer to the wall
and repeats the experiment. What is the new frequency
of clapping?
\hfill\ansref{Ch18-B1}

\vspace{0.4cm}

\textbf{B2.} A train sounds its whistle at
$720\ \text{Hz}$ while travelling at
$25\ \text{m\,s}^{-1}$ toward a stationary
station. The speed of sound is $340\ \text{m\,s}^{-1}$.\\[4pt]
(a) Find the frequency heard by a stationary person
on the platform as the train approaches.\\
(b) Find the frequency heard after the train has
passed and is moving away.\\
(c) Find the change in frequency as the train passes.\\
(d) The train is still approaching but now the
person is running away from the station at
$5\ \text{m\,s}^{-1}$. Find the frequency heard.\\
(e) Both the train and the person are moving toward
each other: train at $25\ \text{m\,s}^{-1}$ and
person at $5\ \text{m\,s}^{-1}$. Find the observed
frequency.
\hfill\ansref{Ch18-B2}

\vspace{0.4cm}

\textbf{B3.} In a resonance tube experiment, a
closed pipe is used with a tuning fork of
$440\ \text{Hz}$. Resonances are observed at
$L_1 = 18.6\ \text{cm}$ and $L_2 = 56.8\ \text{cm}$.\\[4pt]
(a) Calculate the wavelength of the sound.\\
(b) Calculate the speed of sound.\\
(c) Calculate the end correction for this pipe.\\
(d) What temperature does your calculated speed
correspond to?\\
(e) Predict the length of the third resonance
position.
\hfill\ansref{Ch18-B3}

\vspace{0.4cm}

\textbf{B4.} A police radar gun emits microwave
radiation at $24.0\ \text{GHz}$ and detects a
Doppler-shifted reflection from an approaching car.
The reflected frequency is $24.000\,48\ \text{GHz}$.
The speed of microwaves is $3.0\times10^8\ \text{m\,s}^{-1}$.\\[4pt]
(a) Calculate the frequency shift ($\Delta f$).\\
(b) For electromagnetic waves reflected from a
moving object, the Doppler formula gives:
$\Delta f/f_0 = 2v_{car}/c$.
Calculate the speed of the car.\\
(c) Convert the car speed to km\,h$^{-1}$.\\
(d) If the speed limit is $100\ \text{km\,h}^{-1}$,
is the driver speeding?\\
(e) State one advantage of using microwaves rather
than sound waves for speed measurement in traffic
enforcement.
\hfill\ansref{Ch18-B4}

\vspace{0.4cm}

\textbf{B5.} A diagnostic ultrasound machine
operates at $5.0\ \text{MHz}$. The speed of
ultrasound in soft tissue is $1540\ \text{m\,s}^{-1}$.\\[4pt]
(a) Calculate the wavelength of the ultrasound
in tissue.\\
(b) What is the smallest feature that can be
resolved by this ultrasound? (Assume resolution
$\approx \lambda/2$.)\\
(c) A pulse is transmitted and an echo is received
$80\ \mu\text{s}$ later from an organ boundary.
Calculate the depth of the boundary.\\
(d) Explain why higher-frequency ultrasound gives
better resolution but poorer penetration.\\
(e) Calculate the frequency of the ultrasound
reflected from a blood vessel in which blood flows
at $0.30\ \text{m\,s}^{-1}$ directly toward the
transducer. (Use $\Delta f/f_0 = 2v_{blood}/v_{sound}$)
\hfill\ansref{Ch18-B5}

\end{exercisebox}
